% GENERATED FILE. Edit canonical lessons, then run npm run generate:books.
% canonical-source-hash: fnv1a64:9ebedef08c3b7b23
% canonical-lessons: UR-C462-aql, UR-C462-dana, UR-C462-nadan, UR-C462-masum, UR-C462-bahakna

\chapter{Sense and Innocence}
\label{ch:ur-sense-and-innocence}
\hlchaptermodality{462}

\begin{chapteropening}
\textbf{By the end of this chapter:} \emph{I can say whether someone acted wisely or naively.}

Complete the last lesson of chapter 462: I can say whether someone acted wisely or naively.
\end{chapteropening}

\section[\texorpdfstring{\ur{عقل} (aql)}{عقل (aql)}]{\ur{عقل} (aql) — sense; wits}

\label{lesson:UR-C462-aql}



\begin{quote}
Before the new one: say the Urdu for moral conduct, then the Urdu for to show respect.
\end{quote}

\subsection*{You'll want to know: \ur{عقل}}
\textbf{\ur{عقل}} — \emph{aql} — \textquotedblleft{}sense; wits\textquotedblright{}.

\begin{quote}
\ur{اس} \ur{لڑکے} \ur{میں} \ur{بڑی} \ur{عقل} \ur{ہے۔}
\end{quote}

\begin{quote}
\emph{is laṛke meṅ baṛī aql hai.}
\end{quote}

\begin{quote}
This boy has a lot of sense.
\end{quote}

Feminine.

\begin{cousinweb}
\textbf{\ur{عقل}} (\emph{aql}) is usually traced, through Persian, to Arabic \emph{ʿaql}, \textquotedblleft{}reason, intelligence\textquotedblright{}.
\end{cousinweb}

\begin{grammarlens}[title={What you've built}]
One more word for describing what you lived and felt.
\end{grammarlens}

\subsection*{Your turn}
\begin{itemize}
\raggedright
  \item \emph{Say it:} \emph{aql}
  \item \emph{Say it:} \emph{aql}, once more
  \item \emph{Say it:} \emph{is laṛke meṅ baṛī aql hai.} — This boy has a lot of sense.
  \item \emph{Recall:} say the Urdu for moral conduct, then the Urdu for to show respect, then say \emph{aql} again
\end{itemize}

\begin{tcolorbox}[breakable,colback=teal!4,colframe=teal!35!black,title={Before you move on}]
What does \ur{عقل} mean? (\textquotedblleft{}Sense; wits\textquotedblright{}.) Say it once more.
\end{tcolorbox}
\section[\texorpdfstring{\ur{دانا} (dānā)}{دانا (dānā)}]{\ur{دانا} (dānā) — wise}

\label{lesson:UR-C462-dana}



\begin{quote}
Before the new one: say the Urdu for to show respect, then the Urdu for wits.
\end{quote}

\subsection*{You'll want to know: \ur{دانا}}
\textbf{\ur{دانا}} — \emph{dānā} — \textquotedblleft{}wise\textquotedblright{}.

\begin{quote}
\ur{دادا} \ur{بڑے} \ur{دانا} \ur{آدمی} \ur{ہیں۔}
\end{quote}

\begin{quote}
\emph{dādā baṛe dānā ādmī haiṅ.}
\end{quote}

\begin{quote}
My grandfather is a very wise man.
\end{quote}

\begin{cousinweb}
\textbf{\ur{دانا}} (\emph{dānā}) is usually traced to Persian \emph{dānā}, \textquotedblleft{}knowing, wise\textquotedblright{}.
\end{cousinweb}

\begin{grammarlens}[title={What you've built}]
One more word for describing what you lived and felt.
\end{grammarlens}

\subsection*{Your turn}
\begin{itemize}
\raggedright
  \item \emph{Say it:} \emph{dānā}
  \item \emph{Say it:} \emph{dānā}, once more
  \item \emph{Say it:} \emph{dādā baṛe dānā ādmī haiṅ.} — My grandfather is a very wise man.
  \item \emph{Recall:} say the Urdu for to show respect, then the Urdu for wits, then say \emph{dānā} again
\end{itemize}

\begin{tcolorbox}[breakable,colback=teal!4,colframe=teal!35!black,title={Before you move on}]
What does \ur{دانا} mean? (\textquotedblleft{}Wise\textquotedblright{}.) Say it once more.
\end{tcolorbox}
\section[\texorpdfstring{\ur{نادان} (nādān)}{نادان (nādān)}]{\ur{نادان} (nādān) — naive}

\label{lesson:UR-C462-nadan}



\begin{quote}
Before the new one: say the Urdu for wits, then the Urdu for wise.
\end{quote}

\subsection*{You'll want to know: \ur{نادان}}
\textbf{\ur{نادان}} — \emph{nādān} — \textquotedblleft{}naive\textquotedblright{}.

\begin{quote}
\ur{وہ} \ur{ابھی} \ur{نادان} \ur{ہے،} \ur{دنیا} \ur{کو} \ur{نہیں} \ur{جانتا۔}
\end{quote}

\begin{quote}
\emph{voh abhī nādān hai, dunyā ko nahīṅ jāntā.}
\end{quote}

\begin{quote}
He is still naive and does not know the world.
\end{quote}

\begin{cousinweb}
\textbf{\ur{نادان}} (\emph{nādān}) is usually traced to Persian \emph{nādān}, \textquotedblleft{}unknowing\textquotedblright{}: \emph{nā-}, \textquotedblleft{}not\textquotedblright{}, and \emph{dān}, \textquotedblleft{}knowing\textquotedblright{}.
\end{cousinweb}

\begin{grammarlens}[title={What you've built}]
One more word for describing what you lived and felt.
\end{grammarlens}

\subsection*{Your turn}
\begin{itemize}
\raggedright
  \item \emph{Say it:} \emph{nādān}
  \item \emph{Say it:} \emph{nādān}, once more
  \item \emph{Say it:} \emph{voh abhī nādān hai, dunyā ko nahīṅ jāntā.} — He is still naive and does not know the world.
  \item \emph{Recall:} say the Urdu for wits, then the Urdu for wise, then say \emph{nādān} again
\end{itemize}

\begin{tcolorbox}[breakable,colback=teal!4,colframe=teal!35!black,title={Before you move on}]
What does \ur{نادان} mean? (\textquotedblleft{}Naive\textquotedblright{}.) Say it once more.
\end{tcolorbox}
\section[\texorpdfstring{\ur{معصوم} (māsūm)}{معصوم (māsūm)}]{\ur{معصوم} (māsūm) — innocent}

\label{lesson:UR-C462-masum}



\begin{quote}
Before the new one: say the Urdu for wise, then the Urdu for naive.
\end{quote}

\subsection*{You'll want to know: \ur{معصوم}}
\textbf{\ur{معصوم}} — \emph{māsūm} — \textquotedblleft{}innocent\textquotedblright{}.

\begin{quote}
\ur{بچے} \ur{کا} \ur{چہرہ} \ur{بڑا} \ur{معصوم} \ur{ہے۔}
\end{quote}

\begin{quote}
\emph{bachche kā chehrā baṛā māsūm hai.}
\end{quote}

\begin{quote}
The child has a very innocent face.
\end{quote}

\begin{cousinweb}
\textbf{\ur{معصوم}} (\emph{māsūm}) is usually traced, through Persian, to Arabic \emph{maʿṣūm}, \textquotedblleft{}kept safe from wrongdoing\textquotedblright{}.
\end{cousinweb}

\begin{grammarlens}[title={What you've built}]
One more word for describing what you lived and felt.
\end{grammarlens}

\subsection*{Your turn}
\begin{itemize}
\raggedright
  \item \emph{Say it:} \emph{māsūm}
  \item \emph{Say it:} \emph{māsūm}, once more
  \item \emph{Say it:} \emph{bachche kā chehrā baṛā māsūm hai.} — The child has a very innocent face.
  \item \emph{Recall:} say the Urdu for wise, then the Urdu for naive, then say \emph{māsūm} again
\end{itemize}

\begin{tcolorbox}[breakable,colback=teal!4,colframe=teal!35!black,title={Before you move on}]
What does \ur{معصوم} mean? (\textquotedblleft{}Innocent\textquotedblright{}.) Say it once more.
\end{tcolorbox}
\section[\texorpdfstring{\ur{بہکنا} (bahaknā)}{بہکنا (bahaknā)}]{\ur{بہکنا} (bahaknā) — to be led astray}

\label{lesson:UR-C462-bahakna}



\begin{quote}
Before the new one: say the Urdu for naive, then the Urdu for innocent.
\end{quote}

\subsection*{You'll want to know: \ur{بہکنا}}
\textbf{\ur{بہکنا}} — \emph{bahaknā} — \textquotedblleft{}to be led astray\textquotedblright{}.

\begin{quote}
\ur{وہ} \ur{غلط} \ur{دوستوں} \ur{کے} \ur{ساتھ} \ur{بہک} \ur{گیا۔}
\end{quote}

\begin{quote}
\emph{voh ghalat dostoṅ ke sāth bahak gayā.}
\end{quote}

\begin{quote}
He was led astray by the wrong friends.
\end{quote}

Past (he …): \ur{وہ} \ur{بہک} \ur{گیا}, \emph{voh bahak gayā}.

\begin{cousinweb}
\textbf{\ur{بہکنا}} (\emph{bahaknā}) is inherited: CDIAL lists the Hindi cousin \emph{bahaknā}, \textquotedblleft{}to be misled\textquotedblright{}, under Sanskrit \emph{vyáthatē}, \textquotedblleft{}trembles, goes astray\textquotedblright{} (CDIAL 12164).
\end{cousinweb}

\begin{grammarlens}[title={What you've built}]
One more word for describing what you lived and felt.
\end{grammarlens}

\subsection*{Your turn}
\begin{itemize}
\raggedright
  \item \emph{Say it:} \emph{bahaknā}
  \item \emph{Say it:} \emph{bahaknā}, once more
  \item \emph{Say it:} \emph{voh ghalat dostoṅ ke sāth bahak gayā.} — He was led astray by the wrong friends.
  \item \emph{Recall:} say the Urdu for naive, then the Urdu for innocent, then say \emph{bahaknā} again
\end{itemize}

\begin{tcolorbox}[breakable,colback=teal!4,colframe=teal!35!black,title={Before you move on}]
What does \ur{بہکنا} mean? (\textquotedblleft{}To be led astray\textquotedblright{}.) Say it once more.
\end{tcolorbox}
